\documentclass[11pt]{article}
\usepackage[margin=1in]{geometry}
\usepackage{booktabs}
\usepackage{hyperref}
\usepackage{longtable}
\usepackage{xcolor}
\usepackage{tabularx}

\title{BVBRC-95 Independent Replication Report:\\
Critical Evaluation of Short, Long, and Hybrid Metagenome Assembly for ARG Contextualization\\
(Brown et al., 2021)}
\author{Independent Replication (Ollie / rick-stevens-ai)}
\date{2026}

\begin{document}
\maketitle

\section*{Paper under evaluation}
Brown CL, Keenum IM, Dai D, Zhang L, Vikesland PJ, Pruden A. ``Critical evaluation of short, long, and hybrid assembly for contextual analysis of antibiotic resistance genes in complex environmental metagenomes.'' \textit{Scientific Reports} \textbf{11}:3753 (2021). DOI: \href{https://doi.org/10.1038/s41598-021-83081-8}{10.1038/s41598-021-83081-8}. PMID 33580146. PMC7881036.

\textbf{Data:} NCBI BioProject \href{https://www.ncbi.nlm.nih.gov/bioproject/PRJNA527877}{PRJNA527877} --- 123 SRA runs (5 WWTPs $\times$ 2 sample types $\times$ Illumina + Nanopore + 7 assemblers).

\textbf{Verdict:} \textbf{PARTIAL} (LLM-judge confidence 0.78)

\section{Paper summary}
The authors evaluated 7 metagenome assemblers (IDBA-UD, MEGAHIT, metaSpades, Canu, Flye, HybridSpades, OPERA-MS) on 10 wastewater metagenomes from 5 international WWTPs sampled with matched Illumina short-read and Oxford Nanopore MinION long-read sequencing. The central research question: how does the choice of sequencing technology and assembler affect the ability to place antibiotic resistance genes (ARGs) in genomic context --- i.e., co-localized with mobile genetic elements (MGEs) and taxonomic hosts on the same contig.

\section{Testable claims}
\begin{tabularx}{\linewidth}{@{}lXlll@{}}
\toprule
ID & Statement & Type & Testable? & Tested? \\
\midrule
C1 & Short-read and hybrid assemblies produce SIMILAR ARG contextualization patterns & Comparative & Yes & Yes \\
C2 & Long Nanopore (raw or assembled) produces DISTINCT ARG contextualization patterns & Comparative & Yes & Yes \\
C3 & Hybrid assembly recovers longer ARG-containing contigs than short-read alone & Quantitative & Yes & Partial \\
C4 & Coverage regime affects chimerism/inversions (spike-in) & Quantitative & Yes & No \\
C5 & Long-read alone recovers FEWER ARGs due to Nanopore error rate & Quantitative & Yes & Yes \\
\bottomrule
\end{tabularx}

\section{Method (independent rerun)}
\begin{enumerate}
\item \textbf{Metadata \& data availability.} Queried NCBI eutils and ENA \texttt{filereport} for PRJNA527877. Confirmed 123 accessions (raw Illumina, raw Nanopore, and pre-computed assemblies for all 5 WWTPs $\times$ 2 sample types $\times$ 7 assemblers, $\sim$153 Gbp raw). Full de-novo re-assembly was infeasible on this budget; instead, we verified the authors' \textit{published assemblies} and re-annotated ARGs with an independent modern tool.
\item \textbf{Sample and assembler selection.} Chose the USA-1-influent sample (the paper's own representative example) and downloaded all 7 pre-computed assemblies from ENA (Megahit SRR12664619, metaSpades SRR13105837, IDBA-UD SRR12664620, HybridSpades SRR12664586, Canu SRR12664608, Flye SRR12664575, OPERA-MS SRR12664597).
\item \textbf{Assembly stats.} For each assembly, computed contig count, total bp, max length, median length, N50, and counts $\geq$1/5/10/50/100 kb.
\item \textbf{ARG annotation.} Filtered each assembly to contigs $\geq$1 kb and ran NCBI AMRFinder+ v3.12.8 with the 2024-07-22 CARD-derived AMR reference DB (\texttt{--plus} extended set, tblastn/blastx). AMRFinder+ is a stricter, more curated caller than the paper's original Diamond-vs-CARD/ACLAME/PATRIC; absolute counts differ but relative patterns across assemblers are the comparable quantity.
\item \textbf{Cross-assembler comparison.} Per-assembler: \texttt{n\_arg\_hits}, \texttt{n\_unique\_arg\_symbols}, \texttt{n\_arg\_carrying\_contigs} and length distribution. Pairwise Jaccard similarity of ARG-symbol sets computed for all 21 assembler pairs, then aggregated within/between the three categories (short, long, hybrid).
\item \textbf{LLM-judge verdict.} Paper claims + replication result table were scored by Argo \texttt{argo:gpt-5.2} (per hard rule: no regex/rule-based scoring).
\end{enumerate}

\textbf{Compute:} \texttt{uicgpu} (8$\times$A100, 255 cores, 2 TB RAM) at \texttt{/data/stevens/BVBRC-95/}. AMRFinder step used 24--48 threads per job. Total wall time for the 7 ARG annotations: $\sim$3 min.

\section{Results vs paper}

\subsection{Assembly N50 pattern (paper Table 2: hybrid + long $>$ short)}
\begin{tabular}{lrrrrr}
\toprule
Assembler & Type & N50 (bp) & Max contig (bp) & Contigs $\geq$10 kb & Contigs $\geq$50 kb \\
\midrule
Megahit         & short  &     469 &  12{,}505 &   1 &   0 \\
metaSpades      & short  &     372 &  63{,}767 & 356 &   3 \\
IDBA-UD         & short  &     907 &  79{,}408 & 556 &   1 \\
HybridSpades    & hybrid &     430 & 116{,}044 & 794 &   6 \\
OPERA-MS        & hybrid &     544 & \textbf{311{,}842} & 430 &  27 \\
Canu            & long   & 19{,}298 & 231{,}420 & 247 &  46 \\
Flye            & long   & \textbf{45{,}101} & 363{,}209 & 743 & \textbf{140} \\
\bottomrule
\end{tabular}

Pattern matches the paper qualitatively: long-read assemblies dominate N50; hybrid produces longer \textit{maximum} contigs than short-read; short-read has the highest contig count.

\subsection{ARG annotation across assemblers (contigs $\geq$1 kb)}
\begin{tabular}{lrrrrrr}
\toprule
Assembler & ARG hits & Unique syms & ARG contigs & Med ARG-contig & Max ARG-contig & ARG-contigs $\geq$10 kb \\
\midrule
Megahit        & 31 & 30 & 22 & 1{,}653  &  3{,}469 & 0 \\
metaSpades     & 77 & 73 & 55 & 1{,}720  & 10{,}941 & 1 \\
IDBA-UD        & 78 & 74 & 52 & 2{,}306  & 24{,}974 & 2 \\
HybridSpades   & \textbf{79} & \textbf{76} & 56 & 1{,}976 & 14{,}025 & 4 \\
OPERA-MS       & 35 & 33 & 24 & 2{,}321 & \textbf{311{,}842} & 5 \\
Canu           & \textbf{1} & \textbf{1} & 1 & 7{,}062 & 7{,}062 & 0 \\
Flye           & 13 & 13 &  8 & \textbf{25{,}454} & 52{,}465 & 6 \\
\bottomrule
\end{tabular}

\subsection{Cross-assembler ARG-symbol Jaccard by category}
\begin{tabular}{lr}
\toprule
Comparison & Mean pairwise Jaccard \\
\midrule
Within short-read       & 0.498 \\
Within long-read        & 0.077 \\
Within hybrid           & 0.397 \\
\textbf{Short vs Hybrid} & \textbf{0.610} $\leftarrow$ highest cross-category \\
Short vs Long           & 0.095 \\
Hybrid vs Long          & 0.099 \\
\bottomrule
\end{tabular}

\subsection{Claim-by-claim result}
\begin{description}
\item[C1 REPRODUCED.] Short-read and hybrid had the highest cross-category ARG-symbol Jaccard (0.610), well above short-vs-long (0.095) or hybrid-vs-long (0.099) --- matches the paper's abstract statement.
\item[C2 REPRODUCED.] Long-read-only assemblies produced ARG-symbol sets both internally inconsistent (within-long Jaccard 0.077) and highly distinct from short/hybrid ($\sim$0.10).
\item[C3 PARTIAL.] Hybrid produces more long ARG-carrying contigs than short-read (HybridSpades 4$\times$ $\geq$10 kb; OPERA-MS 5$\times$ $\geq$10 kb + 1$\times$ $\geq$50 kb at 311 kb; short-read only 0--2 $\geq$10 kb). Explicit MGE/taxonomy co-carriage (paper's MetaCompare pipeline) not rerun.
\item[C4 NOT-TESTED.] Requires the in-silico \textit{M. hydrocarbonoclasticus} spike-in experiment; out of scope.
\item[C5 REPRODUCED.] Canu detected only 1 ARG, Flye 13, vs 31--79 for short/hybrid --- a $>$2$\times$ to 75$\times$ reduction, consistent with the paper's Nanopore-only depletion finding.
\end{description}

\section{Verdict: PARTIAL}
Three of the paper's five stated claims (C1, C2, C5) were independently reproduced on real author-deposited data using a stricter, independent ARG caller (NCBI AMRFinder+) with quantitatively strong agreement (Jaccard analysis, ARG-count distributions). One claim (C3) was partially reproduced (contig-length dimension confirmed; explicit MGE co-carriage not scored). One claim (C4) was not tested. The N=1 (single sample) scope prevents a ``REPLICATED'' verdict; but the direction and magnitude of the effect clearly hold. This is a \textbf{PARTIAL replication with high confidence in the reproduced claims}.

\section{Genuine Critique}
This section is a candid, non-defensive appraisal of both the underlying paper and this replication attempt. It is meant to inform readers about what is and is not supported by the evidence on disk.

\subsection*{Critique of the paper (Brown et al., 2021)}
\begin{itemize}
\item \textbf{Nanopore-era-dependence.} The paper's chief negative finding --- that long-read-only assemblies collapse ARG ORFs due to indel error --- is tightly coupled to 2019--2020-vintage MinION R9.4.1 chemistry and pre-Bonito basecallers (Guppy $\sim$3.x). Post-2022 Kit 14 / R10.4.1 + Dorado Sup basecalling has substantially closed the ORF-integrity gap, and the paper's cross-technology conclusions should be re-derived on modern data before being used as general policy for ARG surveillance. This is a genuine external-validity concern, not merely a caveat.
\item \textbf{ARG caller choice conflates method with technology.} The paper's Diamond-vs-CARD/ACLAME/PATRIC pipeline gives generous hits on short high-identity ORFs, which happen to be where short-read assemblies excel; the finding that ``long-read has fewer ARGs'' is partly a caller-artifact of homology-search on error-ridden sequences. Any long-read-first ARG pipeline should include indel-tolerant callers (e.g., HMMER on translated 6-frame windows, or graph-based approaches) before concluding long-read is inferior.
\item \textbf{MGE/host co-carriage evidence is contig-anecdotal.} Paper Figures 4--5 rest on visual inspection of a small number of representative contigs. Statistical support for the co-carriage claim is thin --- there is no formal test that hybrid enriches for ARG--MGE proximity relative to a null model of contig-length matched controls.
\item \textbf{Sample scope is small} (5 WWTPs, 10 metagenomes). Generalization to soil/sediment/drinking-water/one-health matrices is asserted, not demonstrated.
\end{itemize}

\subsection*{Critique of this replication}
\begin{itemize}
\item \textbf{N=1 sample.} Only USA-1-influent was analyzed. Any per-sample noise could dominate; the concordant Jaccard result is encouraging but not a population-level replication. A full 10-sample rerun is the honest next step.
\item \textbf{Different ARG caller.} Switching from Diamond-vs-CARD to AMRFinder+ was necessary and defensible but introduces a confound: raw ARG counts are not directly comparable to the paper's numbers. We claim the \textit{cross-assembler pattern} is caller-invariant, but this was not formally shown by running both callers side-by-side.
\item \textbf{Pre-computed assemblies used.} We annotated the authors' deposited assemblies rather than re-assembling from raw reads. This means any bug in the authors' assembly step propagates silently into our replication. A stronger rerun would re-execute at least one assembler (e.g., Flye 2.9.x on modern basecalled data) from raw reads.
\item \textbf{MGE co-carriage not scored.} The paper's central biological claim --- that assembly choice affects ARG-\textit{context} (not just ARG counts) --- is only partially addressed. Contig length is a proxy; explicit MGE and taxonomy annotation on ARG contigs was not run.
\item \textbf{No de-novo spike-in.} C4 (coverage $\to$ chimerism) is untested, so the paper's spike-in argument stands or falls independently of this replication.
\item \textbf{Single ARG DB snapshot} (2024-07-22). Later DB versions could shift some symbols; we did not do a sensitivity analysis across DB versions.
\end{itemize}

\subsection*{What would move the verdict from PARTIAL to REPLICATED}
Running all 10 samples through the same pipeline; running two ARG callers in parallel (AMRFinder+ and Diamond-vs-CARD, as in the paper) to formally show caller-invariance; adding MetaCompare-style MGE co-carriage analysis; and re-basecalling the raw Nanopore data with Dorado Sup and re-assembling with modern Flye 2.9.x to check whether the C5 depletion is chemistry-era-specific.

\section{Evidence}
\begin{itemize}
\item \texttt{report/evidence/assembly\_stats.jsonl} --- raw assembly statistics per assembler
\item \texttt{report/evidence/summary.json} --- combined assembly + AMR summary
\item \texttt{report/evidence/arg\_symbols\_by\_assembler.json} --- full ARG symbol sets per assembler
\item \texttt{report/evidence/*.1kb.amr.tsv} --- raw AMRFinder+ output for all 7 assemblers
\item \texttt{report/evidence/analyze\_amr.sh}, \texttt{report/evidence/filter\_and\_amr.sh} --- analysis pipelines
\item \texttt{report/evidence/analysis\_output.txt} --- captured stdout
\item \texttt{report/evidence/llm\_judge.json} --- LLM-judge verdict (argo:gpt-5.2)
\end{itemize}

\section{Tool versions}
NCBI AMRFinder+ v3.12.8 with database 2024-07-22.1; Python 3.8.10; curl (data download from ENA); data via ENA FTP mirror of SRA.

\section{Limitations}
Single-sample scope, different ARG caller, MGE co-carriage not run, in-silico spike-in not attempted --- see Genuine Critique for the full accounting.

\end{document}
